\documentclass[10pt,a4paper]{amsart}
\usepackage[margin=19mm]{geometry}
\usepackage{amsmath,amssymb,mathtools}
\usepackage[scr=rsfs,cal=euler]{mathalfa}
\usepackage{xfrac}
\usepackage[unicode,hidelinks,bookmarks=false]{hyperref}
\newcommand{\ie}{i.e.\ }
\newcommand{\cf}{cf.\ }
\newcommand{\resp}{respectively\ }
\newcommand{\Subsection}{\subsection}
\providecommand{\nobreakdash}{\nobreak\hbox{-}}
\setlength{\emergencystretch}{2em}
\multlinegap0pt
\usepackage{amssymb}
\usepackage{bm}
\usepackage{xcolor}
\usepackage{dsfont}
\usepackage{tikz}
\newtheorem{lem}{Lemma}
\newtheorem{prop}{Proposition}
\def \Z{\mathbb{Z}}
\newcommand{\bN}{\mathbb{N}}
\newcommand{\bZ}{\mathbb{Z}}
\newcommand{\beq}{\begin{equation}}
\newcommand{\eeq}{\end{equation}}
\newcommand{\gge}{\geqslant}
\newcommand{\scrX}{{\mathscr{X}}}
\newcommand{\tl}{\tilde}
\title{A Drinfeld type presentation of twisted Yangians\\ of quasi-split type}
\author{Kang Lu, Weinan Zhang}
\date{Source: arXiv:2408.06981v2 (CC BY 4.0)}
\begin{document}
\maketitle

\noindent\emph{Source and licence.} A Drinfeld type presentation of twisted Yangians\\ of quasi-split type. Version arXiv:2408.06981v2 (CC BY 4.0), distributed by arXiv under the Creative Commons Attribution 4.0 International licence, http://creativecommons.org/licenses/by/4.0/, as recorded in the arXiv licence metadata for this exact version and shown on the abstract page https://arxiv.org/abs/2408.06981v2. Full licence notice: \texttt{licenses/arxiv-2408.06981-CC-BY-4.0.txt}.\par\smallskip

\noindent\textit{Editorial scope.} Counted unit: the article's prop X=3prop (source0.tex lines 2777--2786). The article's complete original proof of it is delivered with it (source0.tex lines 2787--2845, one \texttt{proof} environment of 59 lines). No mathematics is written, repaired or re-derived: the statement and the proof are the article's own lines, in the article's own notation, and no retained line is substituted or re-flowed.  Retained uncounted as the source-local context the proof needs: lines 1854--1860; lines 1989--2001; lines 2661--2670.  Nothing was omitted from any retained span, so the package carries no omission record.  No retained line contains a citation, so the wrapper carries no bibliography and no bibliography entry.  No retained line contains a figure environment, an image-inclusion command or drawing code, so no image is delivered and the package declares no asset dependency.  Editorial formatting only, in addition: an AMS \texttt{amsart} preamble that restates the article's own declarations and macros used by the retained lines, namely the environments \texttt{lem}, \texttt{prop} on separate counters, together with the standard packages those lines need.  The retained environments print in this document as Lemma 1, Proposition 1 in order of appearance, whereas the article prints the same items with the numbering of its own full document.  The complete native source of the article as distributed in the arXiv e-print for this exact version is bundled unchanged as \texttt{sources/163736/source0.tex}.
\begin{lem}\label{efaiii}
In the algebra $\scrX_N$, we have $c(u)\tl d_{i'}(u)=d_{i}(-u)$,
\beq\label{efd}
\tl d_i(u)d_{i+1}(u)=\tl d_{i'-1}(-u)d_{i'}(-u),\quad  \tl e_{ij}(u)=f_{i'j'}(-u), \quad  \tl f_{ji}(u)=e_{j'i'}(-u).
\eeq
In particular, we have $e_{i}(u)=-f_{\tau i}(-u)$.
\end{lem}

It follows from Lemma \ref{efaiii} that
\begin{align}
\label{i=tauih}
&h_{\tau i}(u)=h_{i}(-u),\\  
&b_{i}(u)=\begin{cases}
-\sqrt{-1}\,e_{\tau i}(-u+\tfrac{n- \tau i}{2}),&\text{ if }N=2n.\\
-\sqrt{-1}\,e_{\tau i}(-u+\tfrac{n- \tau i}{2}+\frac{1}{4}), &\text{ if }N=2n+1.
\end{cases}.\label{fi=tauie}
\end{align}
In particular, if $N=2n$ is even, then $h_n(u)$ is an even series in $u^{-1}$. Moreover, under the anti-automorphism $\eta$, we have
\beq\label{etahb}
\eta (h_i(u))=h_i(u),\qquad \eta(b_{i}(u))=-b_{\tau i}(-u).
\eeq

\begin{align}
[d_i(u),d_j(v)]&=0,\label{ddb3}\\
[d_1(u),e_1(v)]&=\frac{1}{u-v}d_1(u)(e_1(v)-e_1(u)),\label{d1e1b3}\\
[d_1(u),f_1(v)]&=\frac{1}{u-v}(f_1(u)-f_1(v))d_1(u),\label{d1f1b3}\\
[e_1(u),e_1(v)]&=\frac{1}{u-v}(e_1(u)-e_1(v))^2,\label{e1e1b3}\\
[f_1(u),f_1(v)]&=-\frac{1}{u-v}(f_1(u)-f_1(v))^2,\label{f1f1b3}\\
[e_1(u),f_1(v)]&=\frac{1}{u-v}\big(\tl d_1(u)d_2(u)-\tl d_1(v)d_2(v)\big)+\frac{1}{u+v}\big(e_{13}(u)+e_1(u)f_1(v)+f_{31}(v)\big),\label{e1f1b3}\\
[d_2(u),e_1(v)]&=\frac{1}{u-v}d_2(u)(e_1(u)-e_1(v))-\frac{1}{u+v}(e_1(v)+f_2(u))d_2(u),\label{d2e1b3}\\
[d_2(u),f_1(v)]&=\frac{1}{u-v}(f_1(v)-f_1(u))d_2(u)+\frac{1}{u+v}d_2(u)(f_1(v)+e_2(u)).\label{d2f1b3}
\end{align}

\begin{prop}\label{X=3prop}
For $i,j\in\{1,2\}$ and $r,s\in\bN$, we have 
\begin{align}
&[h_{i,r},h_{j,s}]=0,\qquad h_{i,r}=(-1)^{r+1}h_{\tau i,r},\label{hhb3}\\
&[b_{i,r+1},b_{j,s}]  - [b_{i,r },b_{j,s+1 }]  =\frac{c_{ij}}{2}\{b_{i,r },b_{j,s}\}-2\delta_{i,\tau j} (-1)^{r}h_{\tau i,r+s+1 },\label{bbb3}\\
&[h_{i,r+2},b_{j,s}] - [h_{i,r},b_{j,s+2}]=\frac{c_{ij}-c_{i,\tau j}}{2}\{h_{i,r+1},b_{j,s}\}\notag\\
&\qquad\qquad\qquad \qquad \qquad \ \  +\frac{c_{ij}+c_{i,\tau j}}{2} \{h_{i,r},b_{j,s+1}\}+\frac{c_{ij}c_{\tau i,j}}{4} [h_{i,r}, b_{j,s}].\label{hbb3}
\end{align}
Here we allow $r\in\bZ$ in the relation \eqref{hbb3} with $h_{i,-1}=1$ and $h_{i,r}=0$ for $r<-1$.
\end{prop}

\begin{proof}
The relation \eqref{hhb3} is clear from \eqref{i=tauih}, \eqref{ddb3}, and the definition of $h_{i,r}$.

The relation \eqref{bbb3} for $i=j$ is obvious from \eqref{e1e1b3} and \eqref{f1f1b3}. Here for \eqref{e1e1b3}, we apply \eqref{fi=tauie} to get the relation \eqref{bbb3} for $i=j=2$. Then we consider the case $(i,j)=(1,2)$. It follows from \eqref{e1f1b3} and $e_1(u)=-f_2(-u)$ (by Lemma \ref{efaiii}) that
\begin{align*}
(u-v)&[f_1(u),f_2(v)]\\
&=\frac{u-v}{u+v}\big(\tl d_1(u)d_2(u)-\tl d_1(-v)d_2(-v)\big)+e_{13}(-v)-f_2(v)f_1(u)+f_{31}(u).
\end{align*}
Using $\tl d_2(v)d_3(v)=\tl d_1(-v)d_2(-v)$ from Lemma \ref{efaiii} and rewriting it in terms of $b_1(u)$ and $b_2(v)$, we have
\beq\label{mm-00}
\begin{split}
\big(u-v+\tfrac12\big)[b_1(u),b_2(v)]
=\frac{u-v+\frac12}{u+v}\Big(\tl d_2(v-\tfrac14)d_3(v-\tfrac14)-\tl d_1(u+\tfrac14)d_2(u+\tfrac14)\Big)\\-e_{13}(-v+\tfrac14)-b_2(v)b_1(u)-f_{31}(u+\tfrac14).
\end{split}
\eeq
Expanding the RHS in the region $|u|\gg |v|$, then the term
\beq\label{996}
\frac{u-v+\frac12}{u+v} \tl d_1(u+\tfrac14)d_2(u+\tfrac14)=\Big(1-2\big(v-\tfrac14\big)\sum_{k\gge 0}(-1)^k\frac{v^k}{u^{k+1}}\Big)\tl d_1(u+\tfrac14)d_2(u+\tfrac14)
\eeq
does not contribute to the coefficients of $u^{-r-1}v^{-s-1}$ with $r\in\bZ$ and $s\in \bN$. For the relation \eqref{bbb3}, we shall only consider the coefficients of $u^{-r-1}v^{-s-1}$ with $r\in\bN$ and $s\in \bN$. Hence we can drop terms like $\tl d_1(u+\tfrac14)d_2(u+\tfrac14)$, $f_{31}(u+\tfrac14)$, $e_{13}(-v+\tfrac14)$ in \eqref{mm-00} and this gives rise to
\[
\big(u-v+\tfrac12\big)[b_1(u),b_2(v)]\simeq \frac{u-v+\frac12}{u+v} \tl d_2(v-\tfrac14)d_3(v-\tfrac14)-b_2(v)b_1(u).
\]
Thus we have
\[
(u-v)[b_1(u),b_2(v)]\simeq\Big(1-\frac{2( v-\tfrac14)}{u+v}\Big)\tl d_2(v-\tfrac14)d_3(v-\tfrac14)-\frac{1}{2}\big\{b_{1}(u),b_2(v)\big\},
\]
which implies further
\[
(u-v)[b_1(u),b_2(v)]\simeq -\frac{1}{2}\big\{b_1(u),b_2(v)\big\}-\frac{2v}{u+v} h_2(v).
\]
By taking the coefficients of $u^{-r-1}v^{-s-1}$ for $r,s\in\bN$, we obtain the relation \eqref{bbb3} for $(i,j)=(1,2)$.

We consider the relation \eqref{hbb3}. Note that $h_1(u)=h_2(-u)$ by \eqref{i=tauih} and the anti-automorphism $\eta$ sends $b_1(u)$ to $-b_2(-u)$; see \eqref{etahb}. It suffices to consider the case $(i,j)=(1,1)$. It follows from \eqref{d1f1b3} and \eqref{d2f1b3} that
\beq\label{1helper}
\begin{split}
[\tl d_1(u)&d_2(u),f_1(v)]\\&=\frac{2}{u-v}\tl d_1(u)(f_1(v)-f_1(u))d_2(u)+\frac{1}{u+v}\tl d_1(u)d_2(u)(f_1(v)+e_2(u)).    
\end{split}
\eeq
Thus, we have
\beq\label{2hper}
[\tl d_1(u)d_2(u),f_1(v)]\simeq\frac{2}{u-v}\tl d_1(u)f_1(v)d_2(u)+\frac{1}{u+v}\tl d_1(u)d_2(u)f_1(v),
\eeq
where we dropped terms like $\tl d_1(u)f_1(u)d_2(u)$. By \eqref{d1f1b3}, we find
\[
\frac{1}{u-v}\tl d_1(u)f_1(v)\simeq \frac{1}{u-v-1}f_1(v)\tl d_1(u).
\]
Plugging it into \eqref{2hper}, clearing the denominator, substituting $u\to u+\frac14$ and $v\to v+\frac14$ and multiplying both sides by $\sqrt{-1}\big(1+\frac{1}{4u}\big)$, we obtain
\begin{align*}
\big(u+v+\tfrac{1}{2}\big)&\big(u-v-1\big)\big[h_1(u),b_1(v)\big]\\
\simeq\,&2 \big(u+v+\tfrac{1}{2}\big)b_1(v)h_1(u)+ \big(u-v-1\big)h_1(u)b_1(v),
\end{align*}
which gives rise to
\begin{align*}
\big(u^2-v^2\big)\big[h_1(u),b_1(v)\big]
\simeq \frac{3u+v}{2} \big\{h_1(u),b_1(v)\big\}-\frac{1}{2}[h_1(u),b_1(v)].                               
\end{align*}
Taking the coefficients of $u^{-r-1}v^{-s-1}$  with $r\in\bZ,s\in\bN$, one finds the relation \eqref{hbb3} for $(i,j)=(1,1)$. Here we can make $r\in\Z$ as the terms we dropped give rise to $u^{k}v^{l}$ for $l\gge 0$ if we expand the denominator in the region $|u|\gg |v|$ as exhibited in \eqref{996}.
\end{proof}

\end{document}
